\documentclass[10pt,a4paper]{amsart}
\usepackage[margin=19mm]{geometry}
\usepackage{amsmath,amssymb,mathtools}
\usepackage[scr=rsfs,cal=euler]{mathalfa}
\usepackage{xfrac}
\usepackage[unicode,hidelinks,bookmarks=false]{hyperref}
\newcommand{\ie}{i.e.\ }
\newcommand{\cf}{cf.\ }
\newcommand{\resp}{respectively\ }
\newcommand{\Subsection}{\subsection}
\providecommand{\nobreakdash}{\nobreak\hbox{-}}
\setlength{\emergencystretch}{2em}
\multlinegap0pt
\usepackage{amssymb}
\newtheorem{assumption}{Assumption}
\newtheorem{proposition}{Proposition}
\newtheorem{theorem}{Theorem}
\DeclareMathOperator*{\minimize}{minimize}
\title{\LARGE \bf An Adaptive Multi-Parameter ADMM Algorithm for Embedded MPC}
\author{Alberto Zaupa\textsuperscript{\ddag} and Mikael Johansson\textsuperscript{\ddag}}
\date{Source: arXiv:2609.12992v1 (CC BY 4.0)}
\begin{document}
\maketitle

\noindent\emph{Source and licence.} \LARGE \bf An Adaptive Multi-Parameter ADMM Algorithm for Embedded MPC. Version arXiv:2609.12992v1 (CC BY 4.0), distributed by arXiv under the Creative Commons Attribution 4.0 International licence, http://creativecommons.org/licenses/by/4.0/, as recorded in the arXiv licence metadata for this exact version and shown on the abstract page https://arxiv.org/abs/2609.12992v1. Full licence notice: \texttt{licenses/arxiv-2609.12992-CC-BY-4.0.txt}.\par\smallskip

\noindent\textit{Editorial scope.} Counted unit: the article's theorem thm:superlinear\_convergence (source0.tex lines 400--407). The article's complete original proof of it is delivered with it (source0.tex lines 409--573, one \texttt{proof} environment of 165 lines). No mathematics is written, repaired or re-derived: the statement and the proof are the article's own lines, in the article's own notation, and no retained line is substituted or re-flowed.  Retained uncounted as the source-local context the proof needs: lines 133--139; lines 141--148; lines 237--240; lines 325--333; lines 337--346; lines 371--382.  Nothing was omitted from any retained span, so the package carries no omission record.  No retained line contains a citation, so the wrapper carries no bibliography and no bibliography entry.  No retained line contains a figure environment, an image-inclusion command or drawing code, so no image is delivered and the package declares no asset dependency.  Editorial formatting only, in addition: an AMS \texttt{amsart} preamble that restates the article's own declarations and macros used by the retained lines, namely the environments \texttt{assumption}, \texttt{proposition}, \texttt{theorem} on separate counters, together with the standard packages those lines need.  The retained environments print in this document as Assumption 1, Proposition 1, Theorem 1 in order of appearance, whereas the article prints the same items with the numbering of its own full document.  The complete native source of the article as distributed in the arXiv e-print for this exact version is bundled unchanged as \texttt{sources/210168/source0.tex}.
\begin{equation}
\label{eq:admm_mpc}
\begin{aligned}
    \minimize_{w,z} &\quad f(w) + \mathcal{I}_{[a,b]}(z) \\
    \text{subject to} &\quad Gw = z
\end{aligned}
\end{equation}

\begin{equation}
\label{eq:admm_iters}
\begin{aligned}
w^{(k+1)} &= \arg \min_w f(w) + \frac{1}{2} \| Gw - z^{(k)} + \rho^{-1} y^{(k)}\|_{\rho}^2 \\
z^{(k+1)} &= \Pi_{[a,b]} \left(Gw^{(k+1)} + \rho^{-1} y^{(k)} \right) \\
y^{(k+1)} &= y^{(k)} + \rho\, (G w^{(k+1)} - z^{(k+1)})
\end{aligned}
\end{equation}

\begin{equation}
\label{eq:update_rule}
    \rho_i^{(k+1)} =\left( \rho^{(k)}_i \right)^{1-\alpha} \Pi_{[\underline{\rho}, \overline{\rho}]} \left(\tilde{\rho}_i^{(k+1)}\right)^{\alpha} 
\end{equation}

\begin{assumption}[Post-identification regime]
\label{assumption:1}
Let $(w^{\star}, z^{\star}, y^{\star})$ be a primal-dual solution of Problem~\eqref{eq:admm_mpc}, and let $\mathbb{A}$ and $\mathbb{I}$ denote the active and inactive coordinates at the solution. We assume that there exists $\bar{K} \in \mathbb{N}$ such that for all $k \geq \bar{K}$, the ADMM iterates generated by~\eqref{eq:admm_iters} satisfy
\begin{equation}
\label{eq:assumtion_1}
\text{sign}(y_i^{(k)}) = \text{sign}(y_i^{\star}) \quad \forall i =1,2,\dots, m
\end{equation}
where we define $\text{sign}(0) = 0$.
\end{assumption}

\begin{proposition}
\label{prop:assumption_1}
Let Assumption~\ref{assumption:1} hold. Then for all $k > \bar{K}$:
\begin{equation}
\begin{aligned}
    z_{\mathbb{A}}^{(k)} &= z_{\mathbb{A}}^{\star} \\
    z_{\mathbb{I}}^{(k)} &= w_{\mathbb{I}}^{(k)}
\end{aligned}
\end{equation}
\end{proposition}

\begin{proposition}
\label{prop:rho_dynamics}
    The proposed update rule~\eqref{eq:update_rule} with $\underline{\rho} = 0$, $\overline{\rho} = +\infty$ and $\epsilon=0$ is such that when Assumption~\ref{assumption:1} holds, $\rho$ updates satisfy:
    \begin{equation}
    \label{eq:assumption_2}
    \begin{aligned}
        \rho_i^{(k+1)} &\ge \beta\, \rho_i^{(k)} \quad \forall i \in \mathbb{A} \\
        \rho_i^{(k+1)} &\le \gamma\, \rho_i^{(k)} \quad \forall i \in \mathbb{I}
    \end{aligned}
    \end{equation}
for some $\beta > 1$ and $\gamma \in (0, 1)$.
\end{proposition}

\begin{theorem}
\label{thm:superlinear_convergence}
Let $(w^{\star}, z^{\star}, y^{\star})$ be a primal-dual solution of Problem~\eqref{eq:admm_mpc} with $G = I$, and denote by $\mathbb{A}$ and $\mathbb{I}$ the set of active and inactive constraints at the triple. Assume that after iteration $\bar{K}$, the sequence $\{(w^{(k)}, z^{(k)}, y^{(k)})\}$ generated by~\eqref{eq:admm_iters} satisfies Assumption~\ref{assumption:1}, and that $\{\rho^{(k)}\}$ is driven by~\eqref{eq:update_rule} with $\underline{\rho} = 0$, $\overline{\rho} = + \infty$ and $\epsilon=0$. Moreover suppose that there exists a closed ball $\mathcal{N}$ centered at $w^{\star}$ such that $f$ is $\mathcal{C}^2$ in $\mathcal{N}$, $\nabla^2 f(w) \succ 0 \; \forall w \in \mathcal{N}$ and $\{w^{(k)}\}_{k\geq\bar{K}} \in \mathcal{N}$. 
Then there exists $C > 0$ such that:
\begin{equation}
    \|w^{(k+1)} - w^{\star} \| \leq C e^{-\Theta(k^2)} \|w^{(0)} - w^{\star}\| \;.
\end{equation}
\end{theorem}

\begin{proof}
Because we are after an asymptotic convergence rate, in the following we may assume without loss of generality $\bar{K} = 0$.

Notice that given our assumptions on $\mathcal{C}^2$ continuity of $f$ within $\mathcal{N}$, the $w$-update is characterized by the optimality condition:
\begin{equation}
\label{eq:w_optimality}
\nabla f(w^{(k+1)}) + y^{(k)} + \rho^{(k)} (w^{(k+1)} - z^{(k)}) = 0
\end{equation}
and slicing the equation above with respect to $\mathbb{A}$ and $\mathbb{I}$ we get:
\begin{equation}
\label{eq:w_opt_manipulation_1}
\nabla f(w^{(k+1)}) + \begin{bmatrix}
    y_{\mathbb{A}}^{(k)} \\
    0
\end{bmatrix} + \begin{bmatrix}
    \rho_{\mathbb{A}}^{(k)} (w_{\mathbb{A}}^{(k+1)} - z_{\mathbb{A}}^{\star}) \\
    \rho_{\mathbb{I}}^{(k)} (w_{\mathbb{I}}^{(k+1)} - w_{\mathbb{I}}^{(k)})
\end{bmatrix} = 0
\end{equation}
where we used $z^{(k)}_\mathbb{A} = z_\mathbb{A}^{\star}$. $w_\mathbb{I}^{(k)} = z_\mathbb{I}^{(k)}$ and $y_{\mathbb{I}}^{(k)} = 0$, as stated by Proposition~\ref{prop:assumption_1} and Assumption~\ref{assumption:1}. However, by definition of the dual update in~\eqref{eq:admm_iters} and because $z_\mathbb{A}^{(k)} = z_\mathbb{A}^{\star}$ $\forall k$,~\eqref{eq:w_optimality} also implies:
\begin{equation}
\label{eq:w_opt_manipulation_2}
\begin{aligned}
\nabla_{w_{\mathbb{A}}} f(w^{(k+1)}) + y_{\mathbb{A}}^{(k)} + \rho_{\mathbb{A}}^{(k)} (w_{\mathbb{A}}^{(k+1)} - z_{\mathbb{A}}^{(k+1)}) &= 0 \iff \\
\nabla_{w_{\mathbb{A}}} f(w^{(k+1)}) + y_{\mathbb{A}}^{(k+1)} &= 0
\end{aligned}
\end{equation}
where we denote by $\nabla_{w_{\mathbb{A}}}f$ the gradient of $f$ w.r.t. $w_{\mathbb{A}}$. Combining~\eqref{eq:w_opt_manipulation_1} and~\eqref{eq:w_opt_manipulation_2}, observing that $w_{\mathbb{A}}^{\star} = z_{\mathbb{A}}^{\star}$ and defining $e^{(k)} = w^{(k)} - w^{\star}$ we get:
\begin{equation}
\label{eq:gradients_errors}
\nabla f(w^{(k+1)}) + \rho^{(k)} e^{(k+1)} = \begin{bmatrix}
    \nabla_{w_{\mathbb{A}}}f(w^{(k)}) \\
    \rho_{\mathbb{I}}^{(k)} e^{(k)}_{\mathbb{I}}
\end{bmatrix}
\end{equation}
where we added $0 = -w_\mathbb{I}^{*} + w_\mathbb{I}^{*}$ to the $w_\mathbb{I}^{(k+1)} - w_\mathbb{I}^{(k)}$ term in~\eqref{eq:w_opt_manipulation_1}.
Now we rewrite the gradients in \eqref{eq:gradients_errors} in terms of the error sequence. The mean value theorem states:
\begin{equation}
\begin{aligned}
    \nabla f(w^{(k)}) &= \nabla f(w^{\star}) + \\
    &\left( \int_0^1 \nabla^2 f(\tau w^{\star} + (1 - \tau) w^{(k)}) d\tau \right) e^{(k)} \\
    &= \nabla f(w^{\star}) +  H^{(k)} e^{(k)}
\end{aligned}
\end{equation}
where given our assumptions on $\nabla^2f$ we have that there exist $0 < m < M$ such that, for all $k$:
\begin{equation}
\label{eq:H_bounds}
mI \preceq H^{(k)} \preceq M I
\end{equation}
At this point, we observe that $\nabla_{w_\mathbb{I}}f(w^{\star}) = 0$, which follows immediately from the stationarity condition $\nabla f(w^{\star}) + y^{\star} = 0$, combined with $y_\mathbb{I}^{\star} = 0$. Therefore, defining:
\begin{equation}
    H^{(k)} = \begin{bmatrix}
        H_{\mathbb{A}, \mathbb{A}}^{(k)} & H_{\mathbb{A},\mathbb{I}}^{(k)} \\
        H_{\mathbb{I}, \mathbb{A}}^{(k)} & H_{\mathbb{I}, \mathbb{I}}^{(k)}
    \end{bmatrix}
\end{equation}
we can rewrite Equation~\eqref{eq:gradients_errors} as follows:
\begin{equation}
\begin{aligned}
    \nabla f(w^{{\star}}) + (H^{(k+1)} + \rho^{(k)}) e^{(k+1)} = \\
    \begin{bmatrix} \nabla_{w_{\mathbb{A}}}f(w^{\star}) + 
    H^{(k)}_{\mathbb{A},\mathbb{A}}\, e_{\mathbb{A}}^{(k)} + H_{\mathbb{A}, \mathbb{I}}^{(k)} \,e_{\mathbb{I}}^{(k)} \\
    \rho_{\mathbb{I}}^{(k)} e_{\mathbb{I}}^{(k)}
\end{bmatrix}
\end{aligned}
\end{equation}
where the gradients cancel out, leading to:
\begin{equation}
\label{eq:error_recursion}
(H^{(k+1)} + \rho^{(k)}) e^{(k+1)} = \begin{bmatrix}
    H^{(k)}_{\mathbb{A},\mathbb{A}}\, e_{\mathbb{A}}^{(k)} + H_{\mathbb{A}, \mathbb{I}}^{(k)} \,e_{\mathbb{I}}^{(k)} \\
    \rho_{\mathbb{I}}^{(k)} e_{\mathbb{I}}^{(k)}
\end{bmatrix}
\end{equation}
which is a recursion in the error $e^{(k)}$. We will now manipulate Equation~\eqref{eq:error_recursion} in order to obtain a rate of convergence for $\|e^{(k)}\|$. We start by solving for $e_{\mathbb{I}}^{(k+1)}$:
\begin{equation}
\label{eq:solve_I}
\begin{aligned}
    e_{\mathbb{I}}^{(k+1)} = \left(P^{(k+1)}\right)^{-1}\left( -H_{\mathbb{I}, \mathbb{A}}^{(k+1)} e_{\mathbb{A}}^{(k+1)} + \rho_{\mathbb{I}}^{(k)} e_{\mathbb{I}}^{(k)} \right)
\end{aligned}
\end{equation}
where $P^{(k+1)} = H_{\mathbb{I}, \mathbb{I}}^{(k+1)} + \rho_{\mathbb{I}}^{(k)}$. We then substitute in the system for $e_{\mathbb{A}}^{(k+1)}$:
\begin{equation}
\begin{aligned}
\label{eq:solve_A}
    S^{(k+1)} e_{\mathbb{A}}^{k+1} = B^{(k+1)} e^{(k)}
\end{aligned}
\end{equation}
where:
\begin{equation}
    \begin{aligned}
        S^{(k+1)} &= H_{\mathbb{A},\mathbb{A}}^{(k+1)} + \rho_{\mathbb{A}}^{(k)} - H^{(k+1)}_{\mathbb{A}, \mathbb{I}} (P^{(k+1)})^{-1} H_{\mathbb{I},\mathbb{A}}^{(k+1)} \\
        B^{(k+1)} &= \begin{bmatrix} H_{\mathbb{A},\mathbb{A}}^{(k)} & H_{\mathbb{A}, \mathbb{I}}^{(k)} - H_{\mathbb{A},\mathbb{I}}^{(k+1)}(P^{(k+1)})^{-1} \rho_{\mathbb{I}}^{(k)} \end{bmatrix}
    \end{aligned}
\end{equation}
We now make two important observations. First, notice that because Proposition~\ref{prop:rho_dynamics} states that $\rho_{\mathbb{I}}^{(k)}$ is exponentially decreasing and $\|H^{(k)}\| \leq M$, $B^{(k+1)}$ is bounded in norm, uniformly over $k$.
\\
The second observation is that for some $c_0 > 0$, we have:
\begin{equation}
\label{eq:lambda_min_bound}
    \lambda_{\min}\left(S^{(k+1)}\right) > c_0 \beta^{k} \,.
\end{equation}
This is an consequence of the fact that, since $P^{(k+1)} \succ H_{\mathbb{I},\mathbb{I}}^{(k+1)}$:
\begin{equation}
\begin{aligned}
    H_{\mathbb{A},\mathbb{A}}^{(k+1)} - H^{(k+1)}_{\mathbb{A}, \mathbb{I}} (P^{(k+1)})^{-1} H_{\mathbb{I},\mathbb{A}}^{(k+1)} &\succ \\
    H_{\mathbb{A},\mathbb{A}}^{(k+1)} - H^{(k+1)}_{\mathbb{A}, \mathbb{I}} (H_{\mathbb{I}, \mathbb{I}}^{(k+1)})^{-1} H_{\mathbb{I},\mathbb{A}}^{(k+1)} &\succ 0
\end{aligned}
\end{equation}
and that according to Proposition~\ref{prop:rho_dynamics}:
\begin{equation}
    \lambda_{\min}\left(\rho_\mathbb{A}^{(k)}\right) \ge \beta^k \lambda_{\min}\left(\rho^{(0)}\right) \,.
\end{equation}
Equation~\eqref{eq:lambda_min_bound} then implies:
\begin{equation}
    \|\left(S^{(k+1)}\right)^{-1}\| = \frac{1}{\lambda_{\min}\left(S^{(k+1)}\right)} < \frac{1}{c_0} \left(\frac{1}{\beta}\right)^{k}
\end{equation}

Using these two observations together with Equation~\eqref{eq:solve_A}, we deduce that there exists $c_1 > 0$ such that:
\begin{equation}
\label{eq:e_A_norm_bound}
\|e_{\mathbb{A}}^{(k+1)}\| \leq c_1 \left(\frac{1}{\beta}\right)^{k} \|e^{(k)}\|
\end{equation}
Then from Equation~\eqref{eq:solve_I} we get:
\begin{equation}
\begin{aligned}
        \|e_{\mathbb{I}}^{(k+1)}\| \leq \frac{1}{m} \left( \|H_{\mathbb{I},\mathbb{A}}^{(k+1)}\| \|e_{\mathbb{A}}^{(k+1)}\|
        +  \|\rho_{\mathbb{I}}^{(k)}\|\|e_{\mathbb{I}}^{(k)}\|\right)
\end{aligned}
\end{equation}
where we used the fact that:
\begin{equation}
\begin{aligned}
    \|\left( P^{(k+1)} \right)^{-1}\| &= \frac{1}{\lambda_{\min}\left(H_{\mathbb{I}, \mathbb{I}}^{(k+1)} + \rho_{\mathbb{I}}^{(k)}\right)}\\
    &\leq \frac{1}{\lambda_{\min}\left( H_{\mathbb{I}, \mathbb{I}}^{(k+1)} \right)} \\
    &\leq \frac{1}{m} \,.
\end{aligned}
\end{equation}
Therefore, using $\|H^{(k+1)}_{\mathbb{I}, \mathbb{A}} \| \leq \|H^{(k+1)}\| \leq M$, $\|\rho_{\mathbb{I}}^{(k)}\| \leq \gamma^k \|\rho^0\|$ and the bound for $\|e_{\mathbb{A}}^{(k+1)}\|$ from Equation~\eqref{eq:e_A_norm_bound}, we get:
\begin{equation}
\label{eq:e_I_norm_bound}
\begin{aligned}
    \|e_{\mathbb{I}}^{(k+1)}\| \leq \frac{M}{m} c_1 \left(\frac{1}{\beta}\right)^k \|e^{(k)}\| + \frac{\|\rho^0\|}{m} \gamma^k \|e_{\mathbb{I}}^{(k)}\|
\end{aligned}
\end{equation}
Defining $\bar{\gamma} = \max\{\gamma, \frac{1}{\beta}\} \in (0,1)$, Inequalities~\eqref{eq:e_A_norm_bound} and ~\eqref{eq:e_I_norm_bound} together imply that there exists $c_2>0$ such that:
\begin{equation}
    \label{eq:error_norm_bound}
    \|e^{(k+1)}\| \leq c_2 \bar{\gamma}^k \|e^{(k)}\|
\end{equation}
Unrolling~\eqref{eq:error_norm_bound}, we get:
\begin{equation}
    \begin{aligned}
        \|e^{(k+1)}\| &\leq c_2 \bar{\gamma}^k \|e^{(k)}\| \\
        &\leq c_2^2 \bar{\gamma}^{k} \bar{\gamma}^{k-1} \|e^{(k-1)}\| \\
        &\vdots \\
        &\leq c_2^{k+1} \bar{\gamma}^{\left(\sum_{s=1}^k s\right)} \|e^{(0)}\| \\
        &= c_2^{k+1} \bar{\gamma}^{\frac{(k+1)k}{2}} \|e^{(0)}\| \\
        &\leq\exp \left(\frac{\log \bar{\gamma}}{2}(k+1) k + |\log c_2|(k+1)\right) \|e^{(0)}\| \\
        &= \exp \left( -\Theta\left(k^2\right)\right) \|e^{(0)}\|
    \end{aligned}
\end{equation}
where we used the fact that $\bar{\gamma} \in (0, 1)$. 
The $e^{-\Theta(k^2)}$ bound on $\|w^{(k)} - w^\star\|$ then extends directly to the KKT residuals $\|r_p^{(k)}\|_{\infty}$ and $\|r_d^{(k)}\|_{\infty}$.
\end{proof}

\end{document}
